% Imporditud: tools/import_eksperimendid.py
% Allikas: Eesti füüsikaolümpiaadi eksperimendi ülesannete
%   kogu (Rael Kalda, 2023), ülesanne Ü72, lahendus L72.
\setAuthor{Koit Timpmann}
\setRound{piirkonnavoor}
\setYear{2020}
\setNumber{P E2}
\setDifficulty{2}
\setTopic{vedelike mehaanika}
\setVahendid{Kumminiit (lahti lõigatud rahakummiriba), konksuga koormis, mõõtejoonlaud, klaas veega. Vee tihedus on 1000 kg/m$^3$}
\setPunktid{12}
\setAllikas{Eksperimendi kogu Ü72, lahendus lk 65}
\setLahendusseis{toores}

\prob{Kumminiit}
Määrake keha tihedus.

\hint

\solu
Õhus mõjub ülesriputatud kehale raskusjõud Fo = mg, millest on võimalik arvu-

tada keha mass m = Fo . (1 p) g

Kui keha sukeldada täielikult vette, mõjub kehale üleslükkejõud Fy = $\rho$vgVk. (1 p)

Kui keha on sukeldatud vette, on keha mõju riputusvahendile üleslükkejõu võrra

väiksem raskujõust. Fv = Fo $- \rho$vgVk. (1 p)

Eelmisest seosest saab leida keha ruumala

Vk = Fo $-$ Fv . $\rho$v g

Seega, keha tihedus võrdub

$\rho$ = m = Fo $\rho$v g = $\rho$v Fo Fo Fv (1 p) V g(Fo $-$ Fv) $-$

Jõud saab mõõta kumminiidi pikenemise abil. Kumminiidi pikenemine on võrdeline kumminiiti venitava jõuga F = k(l $-$ $l_0$). (1 p) Mõõdame kumminiidi pikkuse, kui niiti ei ole venitatud $l_0$. (1 p) Mõõdame kumminiidi pikenemise, kui keha ripub kumminiidi otsas õhus. Kehtib seos

Fo = mg = k(lo $-$ $l_0$) (1 p)

Mõõdame kumminiidi pikenemise, kui keha on sukeldatud vette

Fv = Fo $-$ Fy = k(lv $-$ $l_0$) (1 p)

Leiame eelmistest seostest keha keskmise tiheduse

$\rho$ = $\rho$v lo $-$ $l_0$ (1 p) lo $-$ lv

Arvutame mõõtmistulemustest keha keskmise tiheduse --- 1 p. Kordusmõõtmised --- 1 p. Õige tulemus täpsusega $\pm$20 \% --- 1 p.
\probend
